% ============================================================
\section{Teorema 4.3 — Independencia del punto base}
% ============================================================

\begin{frame}{Teorema 4.3}

  \begin{block}{Teorema 4.3}
    Si $X$ es un espacio arco-conexo y $x_0, x_1 \in X$, entonces
    \[
      \pione(X, x_0) \;\cong\; \pione(X, x_1).
    \]
  \end{block}

  \bigskip
  \textbf{Idea central:} cualquier camino $\rho$ de $x_0$ a $x_1$ induce
  un isomorfismo por \emph{conjugación}:
  \[
    [\alpha] \;\longmapsto\; [\bar{\rho} * \alpha * \rho].
  \]

  \bigskip
  \begin{itemize}
    \item El resultado justifica hablar de \emph{el} grupo fundamental
          $\pione(X)$ de un espacio arco-conexo (como grupo abstracto).
    \item \textbf{Cuidado:} el isomorfismo \emph{depende} del camino $\rho$
          elegido; no hay un isomorfismo canónico en general.
  \end{itemize}

\end{frame}

\begin{frame}{Construcción del isomorfismo $P$}

  Sea $\rho : I \to X$ un camino con $\rho(0) = x_0$ y $\rho(1) = x_1$.

  \medskip
  \begin{block}{Definición de $P$}
    \[
      P : \pione(X,x_0) \to \pione(X,x_1),
      \qquad
      P([\alpha]) \;=\; [(\bar{\rho} * \alpha) * \rho].
    \]
  \end{block}

  \medskip
  \textbf{$(\bar{\rho} * \alpha) * \rho$ es lazo en $x_1$:}
  \begin{itemize}
    \item $\bar{\rho}$ va de $x_1$ a $x_0$; $\alpha$ es lazo en $x_0$;
          $\rho$ va de $x_0$ a $x_1$. La concatenación empieza y
          termina en $x_1$. \hfill $\checkmark$
  \end{itemize}

  \medskip
  \textbf{$P$ está bien definida:} si $\alpha \sim_{x_0} \alpha'$
  mediante una homotopía de lazos $F : I \times I \to X$, entonces
  \[
    (\bar{\rho} * F(\cdot, s)) * \rho
  \]
  es una homotopía de lazos de $\bar{\rho} * \alpha * \rho$ a
  $\bar{\rho} * \alpha' * \rho$ en $x_1$. Luego $P([\alpha]) = P([\alpha'])$.
  \hfill $\checkmark$

\end{frame}

\begin{frame}{Ilustración — Conjugación por un camino}

  \begin{center}
  \begin{tikzpicture}[scale=1.1]
    % Espacio X
    \draw[thick, fill=gray!10]
      plot[smooth cycle, tension=0.65]
      coordinates {(-3,0) (-2.5,1.8) (-0.5,2.6) (1.8,2.3) (3,0.6)
                   (2.5,-1.2) (0,-2) (-2,-1.6)};

    % Puntos base
    \filldraw[blue]      (-1.1,-0.3) circle (2.2pt)
      node[below left, blue] {$x_0$};
    \filldraw[verdeTeal] ( 1.3, 0.9) circle (2.2pt)
      node[above right, verdeTeal] {$x_1$};

    % Camino rho (x0 -> x1)
    \draw[blue!60!verdeTeal, thick, ->, >=Stealth, shorten >=3pt]
      (-1.1,-0.3) .. controls (-0.2,0.7) and (0.5,0.5) .. (1.3,0.9)
      node[midway, above=2pt] {$\rho$};

    % Lazo alpha en x0
    \draw[coralOscuro, thick, ->, >=Stealth]
      (-1.1,-0.3)
        .. controls (-2.3,0.1) and (-2.6,1.2) .. (-1.9,1.7)
        .. controls (-1.1,2.2) and (-0.3,1.5) .. (-1.1,-0.3);
    \node[coralOscuro] at (-2.0,1.05) {$\alpha$};

    % Lazo P([alpha]) = rho-bar * alpha * rho en x1
    \draw[ambar, thick, dashed, ->, >=Stealth]
      (1.3,0.9)
        .. controls (0.5,0.5) and (-0.2,0.7) .. (-1.1,-0.3)   % rho-bar
        .. controls (-2.3,0.1) and (-2.6,1.2) .. (-1.9,1.7)   % alpha
        .. controls (-1.1,2.2) and (-0.3,1.5) .. (-1.1,-0.3)  % alpha end
        .. controls (-0.2,0.7) and (0.5,0.5) .. (1.3,0.9);    % rho
    \node[ambar] at (0.1,-1.1) {\small $\bar{\rho}*\alpha*\rho$};

    \node at (2.6,2.0) {$X$};
  \end{tikzpicture}
  \end{center}

  \vspace{-2mm}
  El lazo $\bar{\rho}*\alpha*\rho$ (amarillo, punteado) recorre $\bar{\rho}$
  hasta $x_0$, traza el lazo $\alpha$ (rojo), y regresa por $\rho$ a $x_1$.

\end{frame}

\begin{frame}{$P$ es homomorfismo}

  Sean $[\alpha], [\beta] \in \pione(X, x_0)$.

  \begin{exampleblock}{Cálculo}
    \begin{align*}
      P\!\left([\alpha]\cdot[\beta]\right)
        &= P([\alpha * \beta])
         = [\bar{\rho} * \alpha * \beta * \rho] \\[4pt]
        &= [\bar{\rho} * \alpha * \rho * \bar{\rho} * \beta * \rho]
    \end{align*}
  \end{exampleblock}

  \medskip
  \textbf{Paso clave:} se inserta $\rho * \bar{\rho}$ (que es homotópico
  al lazo constante $c_{x_0}$ por el Lema B) entre $\alpha$ y $\beta$.
  El Lema C garantiza que podemos reasociar libremente, luego:
  \[
    = [\bar{\rho}*\alpha*\rho]\cdot[\bar{\rho}*\beta*\rho]
    = P([\alpha])\cdot P([\beta]).
  \]

  \medskip
  Por tanto $P$ es un homomorfismo de grupos. \hfill $\checkmark$

\end{frame}

\begin{frame}{El inverso $Q$ — $QP = \mathrm{id}$}

  \begin{block}{Definición de $Q$}
    \[
      Q : \pione(X,x_1) \to \pione(X,x_0),
      \qquad
      Q([\sigma]) \;=\; [\rho * \sigma * \bar{\rho}].
    \]
  \end{block}

  \medskip
  \textbf{$QP = \mathrm{id}_{\pione(X,x_0)}$:} Sea $[\alpha] \in \pione(X,x_0)$.

  \begin{align*}
    QP([\alpha])
      &= Q\!\left([\bar{\rho}*\alpha*\rho]\right) \\
      &= [\rho * \bar{\rho} * \alpha * \rho * \bar{\rho}] \\
      &= [\rho*\bar{\rho}] \cdot [\alpha] \cdot [\rho*\bar{\rho}].
  \end{align*}

  Por el Lema B, $\rho * \bar{\rho} \sim_{x_0} c_{x_0}$, y el lazo constante
  es el elemento neutro, así:
  \[
    = [c_{x_0}] \cdot [\alpha] \cdot [c_{x_0}] \;=\; [\alpha]. \hfill \checkmark
  \]

\end{frame}

\begin{frame}{$PQ = \mathrm{id}$ — Conclusión del isomorfismo}

  \textbf{$PQ = \mathrm{id}_{\pione(X,x_1)}$:} Por simetría, intercambiamos
  los papeles de $x_0$ y $x_1$ (y de $\rho$ con $\bar{\rho}$).
  Sea $[\sigma] \in \pione(X,x_1)$:

  \begin{align*}
    PQ([\sigma])
      &= P\!\left([\rho*\sigma*\bar{\rho}]\right) \\
      &= [\bar{\rho}*\rho*\sigma*\bar{\rho}*\rho] \\
      &= [\bar{\rho}*\rho] \cdot [\sigma] \cdot [\bar{\rho}*\rho].
  \end{align*}

  Ahora $\bar{\rho}*\rho \sim_{x_1} c_{x_1}$ por el Lema B, luego:
  \[
    = [c_{x_1}] \cdot [\sigma] \cdot [c_{x_1}] \;=\; [\sigma]. \hfill \checkmark
  \]

  \bigskip
  \begin{block}{Conclusión}
    $Q \circ P = \mathrm{id}$ y $P \circ Q = \mathrm{id}$, por lo que $P$
    es biyectiva. Como además es homomorfismo:
    \[
      P : \pione(X, x_0) \;\xrightarrow{\;\cong\;}\; \pione(X, x_1).
      \hfill \square
    \]
  \end{block}

\end{frame}

\begin{frame}{Observación sobre el punto base}

  \begin{block}{Corolario}
    Para un espacio arco-conexo $X$, todos los grupos $\pione(X,x)$
    son isomorfos. Se puede hablar de \emph{el} grupo fundamental
    $\pione(X)$ como grupo abstracto.
  \end{block}

  \bigskip
  \begin{alertblock}{Advertencia: el isomorfismo no es único}
    \begin{itemize}
      \item Diferentes caminos $\rho$ de $x_0$ a $x_1$ producen, en general,
            \textbf{diferentes isomorfismos} $P_\rho$.
      \item $P_\rho$ y $P_{\rho'}$ difieren por un automorfismo interior de
            $\pione(X,x_1)$ (conjugación por $[\bar{\rho}' * \rho]$).
      \item Al comparar $\pione(X,x_0)$ con $\pione(Y,f(x_0))$ a través de
            $f : X \to Y$, el punto base es \textbf{indispensable}.
    \end{itemize}
  \end{alertblock}

\end{frame}
